\chapter{Conclusion}
The primary objective of this thesis was to design and implement a scalable, on-premise infrastructure capable of gathering, processing, and storing large-scale image datasets for machine learning applications. Based on the analysis of distributed storage solutions and the practical deployment within the Recombee infrastructure, the project successfully addressed the small file problem by utilizing SeaweedFS. This architecture, inspired by Facebook’s Haystack, proved to be highly effective at reducing metadata overhead by grouping small images into larger volumes and maintaining indexes in fast memory. Testing of this storage implementation identified no major problems, and the system functioned as expected. However, it is important to note that while the evaluated benchmarks demonstrate superior performance on small-object workloads, they do not necessarily account for all complex variables or potential issues that may arise under different large-scale real-world workloads.

The implemented data processing pipeline successfully realized a hybrid stream-processing model that combines the low latency of online systems with the high throughput of offline batch systems. By leveraging Kafka as a message broker for inter-stage communication and Kubernetes for orchestration, the system achieved a high degree of modularity and fault tolerance. The use of KEDA for event-driven autoscaling ensured that pipeline components like the Downloader and Embedder can scale horizontally in response to consumer lag, optimizing resource efficiency during periods of low activity while maintaining high throughput during peak loads. The modular approach ensures that specific computational bottlenecks, such as intensive embedding calculations, do not stall the ingestion of new data.

The presented solution architecture is effective for building data processing pipelines with storage nodes optimized for small files. The design focuses primarily on small binary files and automated resource allocation. Its results can be applied to tasks with similar requirements, such as managing e-commerce platforms or storing social media photos. To further share the practical findings of this implementation, the deployment in production use of the work will be documented in a blog post, which will present detailed performance results from the production environment.